Furthermore,
we have @@M0@@, where @@M1@@ runs through the set of nonzero elements of @@M2@@. Consequently, @@M3@@, and:
@@M4@@
by EGA IV 8.2.3 and 8.2.5. Since the schemes considered are noetherian, and the transition morphisms affine and flat, we have by IV 3.2.3:
@@M5@@
Since the ``inverse-image'' homomorphisms on the @@M6@@ are linear with respect to the @@M7@@ (IV (2.12.2)), there exists a nonempty open subset @@M8@@ of @@M9@@ such that in @@M10@@ we have:
@@M11@@
Let @@M12@@ be a closed subscheme of @@M13@@ with support @@M14@@. Then @@M15@@ is surjective by the induction hypothesis, and consequently diagram (1.2.1) shows that:
@@M16@@
so, since @@M17@@, @@M18@@ is surjective.

\sgaStableTarget{sga6.e09.cor.1-3}\textbf{Corollary 1.3.} \uline{Let} @@M19@@ \uline{be a flat morphism of noetherian schemes}. \uline{Suppose that for every point} @@M20@@ \uline{of} @@M21@@, \uline{the homomorphism}:
@@M22@@
\uline{is surjective}. \uline{Then the homomorphism obtained from} @@M23@@:
@@M24@@
\uline{is surjective}.
